\chapter{نتیجه‌گیری}

در پروژهٔ حاضر، چالش افت عملکرد ناشی از نقصان‌های برخورد در حافظه‌های نهان با معماری نگاشت مستقیم مورد بررسی و ارزیابی دقیق قرار گرفت. پدیدهٔ تلاطم که در اثر رقابت آدرس‌های متداخل برای قرارگیری در یک خط فیزیکی مشترک ایجاد می‌شود، می‌تواند ترافیک سنگین و بیهوده‌ای را به سطوح پایین‌تر حافظه تحمیل کند. برای رفع این مشکل، یک پیمانهٔ سخت‌افزاری جدید به نام حافظهٔ نهان قربانی طراحی و در بستر شبیه‌ساز آکیتا پیاده‌سازی شد.

بررسی طراحی و معماری این ماژول نشان داد که اضافه کردن یک بافر بسیار کوچک تمام‌انجمنی (با ظرفیت تنها هشت بلاک) میان حافظهٔ نهان سطح یک و حافظهٔ اصلی، می‌تواند به عنوان یک لایهٔ محافظ در برابر اخراج‌های مکرر عمل کند. پیاده‌سازی مکانیزم جابه‌جایی دوطرفه و حفظ خاصیت انحصاری داده‌ها تضمین کرد که ظرفیت محدود این بافر بهینه‌ترین استفاده را داشته باشد و از ذخیرهٔ داده‌های تکراری در سطوح مختلف جلوگیری شود. افزون بر این، با استفاده از سیاست جایگزینی کمترین استفادهٔ اخیر، بلوک‌های قدیمی‌تر به درستی اخراج شده و در صورت نیاز به حافظهٔ اصلی فرستاده شدند تا داده‌ای از دست نرود. یکی از دستاوردهای مهم در بخش توسعهٔ کد، مدیریت صحیح پیام‌های ناهمگام و جلوگیری از بروز بن‌بست در شبکهٔ ارتباطی شبیه‌ساز بود که به کمک ایجاد صف برای درخواست‌های در انتظار محقق گردید.

نتایج به دست آمده از اجرای شبیه‌ساز بر روی برنامه‌های محک به روشنی اثربخشی این معماری را ثابت کرد. در آزمون‌های طراحی شده که تلاطم شدیدی را به سامانه اعمال می‌کردند، نرخ اصابت در حافظهٔ نهان سطح یک به شدت کاهش یافت؛ اما حافظهٔ نهان قربانی توانست بیش از نیمی از نقصان‌های سطح یک را درون خود مهار کند. این نرخ اصابت بالا در بافر قربانی، باعث کاهش ۵۰ درصدی ترافیک ارسالی به حافظهٔ اصلی شد و در نتیجه، زمان کل اجرای برنامه بیش از ۴۰ درصد بهبود یافت.

تحلیل و مقایسهٔ زمان دسترسی میانگین در دو حالت تئوری و عملی نیز نشان‌دهندهٔ بهبودهای چشمگیر بود. اگرچه محدودیت‌های پنهان در شبکه مانند ترافیک ناشی از عملیات بازنویسی و توقف‌های ساختاری به دلیل پر شدن ثبات درخواست‌های معلق باعث می‌شوند سامانه در عمل کندتر از محاسبات ریاضی رفتار کند، اما استفاده از حافظهٔ قربانی در هر دو حالت توانست تاخیر کلی سامانه را به میزان قابل توجهی کاهش دهد و اثر مخرب تاخیر صد چرخه‌ای حافظهٔ اصلی را خنثی کند.

در پایان، این شبیه‌سازی به طور عملی نشان داد که ایجاد موازنه میان هزینهٔ سخت‌افزاری و بهبود عملکرد چگونه در طراحی معماری کامپیوتر نقش کلیدی ایفا می‌کند. با صرف یک هزینهٔ بسیار ناچیز در سطح سخت‌افزار برای ساخت یک بافر کوچک، سامانه توانست از جریمه‌های سنگین حافظه فرار کرده و پایداری عملکرد خود را در برابر بدترین و پیچیده‌ترین الگوهای دسترسی به حافظه حفظ کند.